\Needspace{5\baselineskip}\textbf{Problems}

\begin{enumerate}[label=\thesection.\arabic*.,ref=\thesection.\arabic*,start=39,series=book-problems]

\item\label{ex:bounded-compact-operator-products} Let $A\in\mathcal{L}(X)$ be compact and let $B\in\mathcal{L}(X)$ be bounded. Show that $AB$ and $BA$ are compact.

\item\label{ex:finite-rank-spectrum-alternative-proof} Prove Lemma~\ref{lem:finite-rank-spectrum-alternative}. Use appropriate results from linear algebra.

\item\label{ex:compact-nonzero-eigenspace-finite-dimensional} Let $A\in\mathcal{L}(X)$ be compact. Show that for $\lambda\ne0$ the solution space of $A\phi=\lambda\phi$ is finite-dimensional.

\item\label{ex:compact-operator-singular-system-expansion} Let $A\in\mathcal{L}(H)$ be compact. Show that there exist orthonormal sets $\{\psi_i\}_{i=1}^N$ and $\{\phi_i\}_{i=1}^N$ and positive real numbers $\{\lambda_i\}_{i=1}^N$ (here $N$ may be finite or infinite) such that
\begin{equation}\label{eq:compact-operator-singular-system-problem}
A(u)=\sum_{i=1}^N\lambda_i(\psi_i,u)\phi_i.
\end{equation}
Hint: $A^*A$ is compact and self-adjoint.

\item\label{ex:compact-operator-norm-preserving-extension} Show that if a compact operator $(\mathcal{D}(A),A)$ from $X$ to $Y$ is extended using the methods defined in case 1 and case 2 of the proof of Theorem~\ref{thm:bounded-operator-norm-preserving-extension}, the extension is also a compact operator.

\item\label{ex:compact-adjoint-hilbert-proof} Let $H$ be a Hilbert space. Prove Theorem~\ref{thm:compact-operator-adjoint-compact} in the case where $X=Y=H$. Hint: Use Theorem~\ref{thm:compact-hilbert-finite-rank-approximation}.

\item\label{ex:relative-compactness-arbitrarily-small-relative-bound} We say that $B$ is compact relative to $A$ if $\mathcal{D}(B)\supseteq\mathcal{D}(A)$ and if $Bx_n$ has a convergent subsequence whenever $x_n\in\mathcal{D}(A)$ and $\|Ax_n\|_Y+\|x_n\|_X$ is bounded. Assume that $A$ is closed, $B$ is closable and that $B$ is compact relative to $A$. Show that $B$ is bounded relative to $A$ and the constant $a$ in Problem~\ref{ex:closed-unbounded-resolvent-analyticity} can be made arbitrarily small. Hint: Try to imitate the proof of Ehrling\textquotesingle s lemma (Lemma~\ref{thm:ehrling-compact-intermediate-norm}).

\item\label{ex:riesz-subspace-separation} Prove the following results due to F. Riesz. Let $S_1$ and $S_2$ be subspaces of a normed linear space. Suppose that $S_1$ is closed and that $S_1$ is a proper subset of $S_2$. Then for every $\theta\in(0,1)$ there is an $x\in S_2$ such that $\|x\|=1$ and
\begin{equation}\label{eq:riesz-subspace-separation-problem}
\|x-y\|\geq\theta\quad\text{for all }y\in S_1.
\end{equation}
Hint: Let $S_2\ni w\notin S_1$, and let $d=\operatorname{dist}(w,S_1)$. Show that there exists $v\in S_1$ such that
\begin{equation}\label{eq:riesz-subspace-separation-distance-hint}
d\leq\|w-v\|\leq\dfrac{d}{\theta}.
\end{equation}


\item\label{ex:infinite-dimensional-identity-noncompact} Show that the unit ball in a normed space $X$ is compact if and only if $X$ is finite-dimensional. Hint: Use Problem~\ref{ex:riesz-subspace-separation}.

\end{enumerate}
